% Barrett modular reduction pipeline — 4-stage CGRA mapping
% Included from 03_openedge.tex via \input
\begin{figure}[t]
\centering
\resizebox{\textwidth}{!}{%
\begin{tikzpicture}[>=Stealth, font=\sffamily]

\def\boxW{4.9}
%\def\boxH{4.6}
\def\boxH{6.0}
\def\gap{0.4}

% \def\barrettbox#1#2#3#4#5#6#7{
%     \begin{scope}[shift={(#1*\boxW + #1*\gap, 0)}]
%     \node[draw=cgragray, thick, rounded corners=3pt, fill=white,
%           minimum width=\boxW cm, minimum height=\boxH cm, anchor=north west]
%           at (0,0) {};
%     \fill[cgralight, rounded corners=3pt] (0.02, -0.02) rectangle (\boxW-0.02, -1.6);
%     \fill[cgralight] (0.02, -1.4) rectangle (\boxW-0.02, -1.6);
%     \draw[cgragray, thick, densely dashed] (0.15, -1.6) -- (\boxW-0.15, -1.6);
%     \node[anchor=north, font=\bfseries\fontsize{8}{9.5}\selectfont, text=cgrablue,
%           text width=\boxW cm - 0.2cm, align=center]
%           at (\boxW/2, -0.15) {#2};
%     \node[anchor=north, font=\fontsize{8}{9.5}\selectfont, align=center,
%           text width=\boxW cm - 0.2cm, text=black]
%           at (\boxW/2, -0.65) {#3};
%     \node[anchor=north, font=\bfseries\fontsize{7.5}{9}\selectfont, text=cgrared,
%           align=center, text width=\boxW cm - 0.2cm]
%           at (\boxW/2, -1.8) {#4};
%     \node[anchor=north, font=\fontsize{7}{8.5}\selectfont, align=center,
%           text width=\boxW cm - 0.4cm, text=black!85]
%           at (\boxW/2, -2.15) {#5};
%     \fill[codebg, rounded corners=2pt, draw=black!10]
%         (0.15, -3.3) rectangle (\boxW-0.15, -4.45);
%     \node[anchor=south east, font=\sffamily\bfseries\fontsize{5.5}{6}\selectfont,
%           text=cgrablue!90!black, fill=cgralight, rounded corners=2pt,
%           inner xsep=4pt, inner ysep=2pt, draw=cgragray!70]
%           at (\boxW+0.2, -4.7) {Map $\rightarrow$ \texttt{#7}};
%     \node[anchor=north west, font=\ttfamily\fontsize{6.5}{8}\selectfont, align=left,
%           text width=\boxW cm - 0.5cm, text=cgragreen!70!black]
%           at (0.25, -3.45) {#6};
%     \coordinate (EAST#1) at (\boxW, -\boxH/2);
%     \coordinate (WEST#1) at (0, -\boxH/2);
%     \end{scope}
% }

\def\barrettbox#1#2#3#4#5#6#7{
    \begin{scope}[shift={(#1*\boxW + #1*\gap, 0)}]
    % 1. Aumentamos la altura visual de la caja si es necesario (ajusta \boxH en tus definiciones)
    \node[draw=cgragray, thick, rounded corners=3pt, fill=white,
          minimum width=\boxW cm, minimum height=\boxH cm, anchor=north west]
          at (0,0) {};
          
    % 2. Ajustamos el sombreado superior para que cubra el nuevo tamaño del título
    \fill[cgralight, rounded corners=3pt] (0.02, -0.02) rectangle (\boxW-0.02, -1.8);
    \fill[cgralight] (0.02, -1.6) rectangle (\boxW-0.02, -1.8);
    \draw[cgragray, thick, densely dashed] (0.15, -1.8) -- (\boxW-0.15, -1.8);
    
    % Título: Bajamos un poco el inicio (-0.15 -> -0.2)
    \node[anchor=north, font=\bfseries\fontsize{10}{12}\selectfont, text=cgrablue,
          text width=\boxW cm - 0.2cm, align=center]
          at (\boxW/2, -0.2) {#2};
          
    % Subtítulo: Bajamos la posición para que no toque al título (-0.65 -> -0.85)
    \node[anchor=north, font=\fontsize{10}{12}\selectfont, align=center,
          text width=\boxW cm - 0.2cm, text=black]
          at (\boxW/2, -0.85) {#3};
          
    % Texto Rojo: Bajamos la posición (-1.8 -> -2.1)
    \node[anchor=north, font=\bfseries\fontsize{9}{11}\selectfont, text=cgrared,
          align=center, text width=\boxW cm - 0.2cm]
          at (\boxW/2, -2.1) {#4};
          
    % Descripción inferior: Bajamos la posición (-2.15 -> -2.55)
    \node[anchor=north, font=\fontsize{8.5}{10}\selectfont, align=center,
          text width=\boxW cm - 0.4cm, text=black!85]
          at (\boxW/2, -2.55) {#5};
          
    % Bloque de Código: Lo desplazamos más abajo para que no choque con la descripción (-3.3 -> -3.8)
    \fill[codebg, rounded corners=2pt, draw=black!10]
        (0.15, -4) rectangle (\boxW-0.15, -6.2);
        
    \node[anchor=north east, font=\sffamily\bfseries\fontsize{7.5}{9}\selectfont,
          text=cgrablue!90!black, fill=cgralight, rounded corners=2pt,
          inner xsep=5pt, inner ysep=3pt, draw=cgragray!70]
          at (\boxW-0.15, -6.2) {Map $\rightarrow$ \texttt{#7}};
          
          
    % Texto del código: Ajustamos acorde al nuevo rectángulo (-3.45 -> -3.95)
    \node[anchor=north west, font=\ttfamily\fontsize{8}{10}\selectfont, align=left,
          text width=\boxW cm - 0.5cm, text=cgragreen!70!black]
          at (0.25, -3.95) {#6};
          
    \coordinate (EAST#1) at (\boxW, -\boxH/2);
    \coordinate (WEST#1) at (0, -\boxH/2);
    \end{scope}
}

\barrettbox{0}
    {Step 1: Base Product}
    {$w \leftarrow x_1 \times x_2$}
    {SPATIAL UNROLLING}
    {16 PEs evaluate all partial products concurrently in one clock bundle.}
    {\textcolor{cgragray}{// subrC (square): 1 bundle}\\[0.4ex]
     bundle \{ at all: R2 = R0 * R1; \}\\[0.4ex]
     \textcolor{cgragray}{// subrB: same, + separate loads}}
    {subrC / subrB}

\barrettbox{1}
    {Steps 2 \& 3: Estimation}
    {$t \leftarrow (w \gg (n{-}1)) \cdot z$\\[0.4ex]$u \leftarrow (t \gg (n{+}1)) \cdot p$}
    {MESH ROUTING}
    {Inter-PE data movement via neighbor registers (RCR, RCB). Shifts as single-bundle ALU ops.}
    {\textcolor{cgragray}{// Quotient estim.: 12 PEs}\\[0.6ex]
     bundle \{ at @0..2,0..3: R2 = R0 * R1; \}}
    {subrA}

\barrettbox{2}
    {Step 4: Mod Subtraction}
    {$(w \bmod 2^{n+1}) - (u \bmod 2^{n+1})$}
    %{EXPLICIT TRUNCATION}
    {EXPL. TRUNCATION}
    {Modulo $2^{n+1}$ via single-bundle AND masking (\texttt{R3 = R3 \& R2}). No division required.}
    {\textcolor{cgragray}{// \dots (carry propagation,}\\[0.6ex]
     \textcolor{cgragray}{// \phantom{\dots }normalization)}}
    {subrA}

\barrettbox{3}
    {Steps 5 \& 6: Correction}
    {$\mathbf{if}\ y > p\ \mathbf{then}\ y \leftarrow y - p$}
    {HW PREDICATION}
    {Branches stall pipelines. Executed via 1-bundle branchless multiplexer.}
    {\textcolor{cgragray}{// Branchless correction:}\\[0.4ex]
     bundle \{ @0,0: R2 = R0 - RCL; \}\\[0.4ex]
     bundle \{ @0,0: BSFA R3\dots \}}
    {subrA}

\foreach \i [evaluate=\i as \next using int(\i+1)] in {0,1,2} {
    \draw[-{Stealth[length=6pt, width=5pt]}, very thick, draw=cgrablue!80!black]
        ([xshift=0.05cm]EAST\i) -- ([xshift=-0.05cm]WEST\next);
}

\end{tikzpicture}%
}
\caption{Barrett modular multiplication mapped onto the CGRA. Each box shows the mathematical operation (top), hardware strategy (middle), and \castm{} code excerpt (bottom).}
\label{fig:barrett-structure}
\end{figure}
